\documentclass[11pt,a4paper]{article}
\usepackage{fontspec,amsmath,amssymb,amsthm,microtype}
\setmainfont{lmroman10-regular.otf}[BoldFont=lmroman10-bold.otf,ItalicFont=lmroman10-italic.otf,BoldItalicFont=lmroman10-bolditalic.otf]
\usepackage{xcolor}
\usepackage[most]{tcolorbox}
\usepackage{mathrsfs}
\providecommand{\lt}{<}
\providecommand{\gt}{>}
\setmonofont{DejaVuSansMono.ttf}[BoldFont=DejaVuSansMono-Bold.ttf,ItalicFont=DejaVuSansMono-Oblique.ttf,BoldItalicFont=DejaVuSansMono-BoldOblique.ttf,Scale=0.86]
\newfontfamily\leanfallback{STIXTwoMath-Regular.otf}[Scale=0.86]
\definecolor{leanrule}{gray}{0.62}
\definecolor{leanback}{gray}{0.962}
\definecolor{leantitle}{gray}{0.22}
\definecolor{leanmuted}{gray}{0.38}
\newtcolorbox{leanbox}[1]{enhanced,breakable,colback=leanback,colframe=leanrule,
  boxrule=0.4pt,arc=0.6mm,left=2.4mm,right=2.4mm,top=1.8mm,bottom=1.8mm,
  before skip=0.8em,after skip=1em,fonttitle=\small\bfseries,coltitle=leantitle,
  colbacktitle=leanback,titlerule=0.3pt,title={#1}}
\newlength{\leancharwidth}
\newenvironment{leancode}{\par\smallskip\begingroup\ttfamily\footnotesize\raggedright
  \settowidth{\leancharwidth}{M}\setlength{\parskip}{0pt}\setlength{\parindent}{0pt}
  \frenchspacing}{\par\endgroup\smallskip}
\newcommand{\leanline}[2]{\hangindent=\dimexpr#1\leancharwidth+4\leancharwidth\relax
  \hangafter=1\noindent\hspace*{#1\leancharwidth}#2\par}
\newcommand{\leansource}[1]{\par{\footnotesize\color{leanmuted}#1}\par}
\newcommand{\leanlabel}[1]{\par\smallskip\noindent{\small\bfseries #1}\par}
\usepackage[margin=25mm]{geometry}
\usepackage[unicode,hidelinks]{hyperref}
\hypersetup{pdfauthor={LangevinLean Project}}
\setlength{\emergencystretch}{3em}
\setlength{\parskip}{0.4em}
\setlength{\parindent}{0pt}
\allowdisplaybreaks[2]
\raggedbottom
\title{\texorpdfstring{Local error from freezing the force}{Local error from freezing the force}}
\author{LangevinLean Project}
\date{}
\begin{document}\maketitle
Draft mathematical exposition; these notes may contain mistakes.


The velocity defect is an explicit exponentially weighted integral of the changing force; the position defect is its primitive. Bounds follow from a Lipschitz force and a kinematic identity, not from a built-in defect hypothesis. The ambient norm may represent a deterministic Banach path or an \(L^{2}\) random-vector norm. A kinetic stationarity theorem supplying the required velocity-moment bound remains separate.



The estimate uses the fact that position is a time integral of velocity. A linear-in-time force variation creates a quadratic velocity defect and a cubic position defect under a shared-noise comparison.


\section{\texorpdfstring{The force-freezing remainder as an integral}{The force-freezing remainder as an integral}}


Set \(D_v(t)=\int_0^te^{-\gamma(t-s)}(g(x_s)-g(x_0))ds\), and \(D_x(t)=\int_0^tD_v(s)ds\). These are exactly the terms left when one compares an exact kinetic equation with a frozen-force interpolation driven by the same additive noise. the velocity defect and the position defect define those vector-valued interval integrals. The definitions impose no desired error estimate. Interpreting them as a numerical-process difference additionally requires deriving the variation-of-constants identity for that process.



The comparison starts from the same phase and uses the same forcing. Every noise term therefore disappears from the error equation. What remains is the effect of replacing the exact moving force by its initial value, so the defect can be represented by an ordinary time integral.

\begin{leanbox}{Lean: \texttt{velocityDefect}, \texttt{positionDefect}}
\leansource{\texttt{velocityDefect} \quad KineticLangevin/\allowbreak{}LocalError.lean, lines 23–24}
\begin{leancode}
\leanline{0}{def\ velocityDefect\ (g\ :\ E\ →\ E)\ (x\ :\ ℝ\ →\ E)\ (γ\ t\ :\ ℝ)\ :\ E\ :=}
\leanline{2}{∫\ s\ in\ 0..t,\ Real.exp\ (-γ\ *\ (t\ -\ s))\ •\ (g\ (x\ s)\ -\ g\ (x\ 0))}
\end{leancode}
\leansource{\texttt{positionDefect} \quad KineticLangevin/\allowbreak{}LocalError.lean, lines 27–28}
\begin{leancode}
\leanline{0}{def\ positionDefect\ (g\ :\ E\ →\ E)\ (x\ :\ ℝ\ →\ E)\ (γ\ t\ :\ ℝ)\ :\ E\ :=}
\leanline{2}{∫\ s\ in\ 0..t,\ velocityDefect\ g\ x\ γ\ s}
\end{leancode}
\leanlabel{How the Lean proof follows the argument}
\texttt{velocity\allowbreak{}Defect} names the mathematical object defined above. Its defining expression appears in the adjacent statement; subsequent identities may unfold this definition directly.




\texttt{position\allowbreak{}Defect} names the mathematical object defined above. Its defining expression appears in the adjacent statement; subsequent identities may unfold this definition directly.
\end{leanbox}

\section{\texorpdfstring{Integrating velocity gives a linear position increment}{Integrating velocity gives a linear position increment}}


Assume \(x_t-x_0=\int_0^tv_sds\), with \(\|v_s\|\le C\) on the horizon. The norm-of-integral bound gives \(\|x_t-x_0\|\le C|t|=Ct\) for nonnegative time. The argument rewrites with the kinematic identity, applies the constant bound on the interval integral, and translates membership in the unoriented integration interval into \(0\le s\le t\le T\). This is the extra position regularity that kinetic Langevin gains from integrating velocity instead of receiving direct Brownian position noise.



If velocity is bounded by \(R\) on the step, then \(\|x_s-x_0\|\le Rs\). This uses only the kinematic relation \(x_s-x_0=\int_0^sv_rdr\). It explains why the force variation is linear in elapsed time even though velocity may be Brownian-driven.

\begin{leanbox}{Lean: \texttt{kinematic\_\allowbreak{}increment\_\allowbreak{}bound}}
\leansource{\texttt{kinematic\_\allowbreak{}increment\_\allowbreak{}bound} \quad KineticLangevin/\allowbreak{}LocalError.lean, lines 32–41}
\begin{leancode}
\leanline{0}{theorem\ kinematic\_increment\_bound\ (x\ v\ :\ ℝ\ →\ E)\ (T\ C\ :\ ℝ)}
\leanline{4}{(hbound\ :\ ∀\ s\ ∈\ Icc\ 0\ T,\ ‖v\ s‖\ ≤\ C)}
\leanline{4}{(hkin\ :\ ∀\ s\ ∈\ Icc\ 0\ T,\ x\ s\ -\ x\ 0\ =\ ∫\ r\ in\ 0..s,\ v\ r)}
\leanline{4}{\{t\ :\ ℝ\}\ (ht\ :\ t\ ∈\ Icc\ 0\ T)\ :\ ‖x\ t\ -\ x\ 0‖\ ≤\ C\ *\ t}
\end{leancode}
\leanlabel{How the Lean proof follows the argument}
\texttt{kinematic\_\allowbreak{}increment\_\allowbreak{}bound} gives the formal counterpart of this step. The main intermediate facts are \texttt{hb} in that order. Named proof applications include \texttt{hbound}; the adjacent proof shows their hypotheses and order of use.
\end{leanbox}

\section{\texorpdfstring{Integrate a linear force increment to obtain the quadratic velocity defect}{Integrate a linear force increment to obtain the quadratic velocity defect}}


Lipschitzness gives \(\|g(x_s)-g(x_0)\|\le MCs\). With \(\gamma\ge0\) and \(s\le t\), the exponential kernel is between zero and one, so the norm of the weighted integrand is still at most \(MCs\). The proof derives this pointwise estimate from the preceding increment theorem, explicit exponential positivity and the norm of scalar multiplication. A continuous scalar majorant provides the required interval integrability. Integrating the bound gives \(\|D_v(t)\|\le\int_0^tMCsds=MCt^2/2\); The argument evaluates the integral and normalizes its coefficients with expansion of the scalar polynomials.



An \(M\)-Lipschitz force gives \(\|g(x_s)-g(x_0)\|\le MRs\). Dropping the exponential kernel only increases the integral when friction is nonnegative. Hence the velocity defect is bounded by \(MR\int_0^h s\,ds=MRh^2/2\).

\begin{leanbox}{Lean: \texttt{velocityDefect\_\allowbreak{}bound}}
\leansource{\texttt{velocityDefect\_\allowbreak{}bound} \quad KineticLangevin/\allowbreak{}LocalError.lean, lines 45–71}
\begin{leancode}
\leanline{0}{theorem\ velocityDefect\_bound\ (g\ :\ E\ →\ E)\ (M\ :\ ℝ≥0)\ (hg\ :\ LipschitzWith\ M\ g)}
\leanline{4}{(x\ v\ :\ ℝ\ →\ E)\ (γ\ T\ C\ :\ ℝ)\ (hγ\ :\ 0\ ≤\ γ)\ (hC\ :\ 0\ ≤\ C)}
\leanline{4}{(hbound\ :\ ∀\ s\ ∈\ Icc\ 0\ T,\ ‖v\ s‖\ ≤\ C)}
\leanline{4}{(hkin\ :\ ∀\ s\ ∈\ Icc\ 0\ T,\ x\ s\ -\ x\ 0\ =\ ∫\ r\ in\ 0..s,\ v\ r)}
\leanline{4}{\{t\ :\ ℝ\}\ (ht\ :\ t\ ∈\ Icc\ 0\ T)\ :}
\leanline{4}{‖velocityDefect\ g\ x\ γ\ t‖\ ≤\ (M\ :\ ℝ)\ *\ C\ *\ t\ \textasciicircum{}\ 2\ /\ 2}
\end{leancode}
\leanlabel{How the Lean proof follows the argument}
\texttt{velocity\allowbreak{}Defect\_\allowbreak{}bound} gives the formal counterpart of this step. The main intermediate facts are \texttt{hb}, \texttt{hint} in that order. Named proof applications include \texttt{Real.\allowbreak{}exp\_\allowbreak{}le\_\allowbreak{}one\_\allowbreak{}iff.\allowbreak{}mpr}; the adjacent proof shows their hypotheses and order of use.
\end{leanbox}

\section{\texorpdfstring{Integrate once more to obtain the cubic position defect}{Integrate once more to obtain the cubic position defect}}


Apply the proved velocity-defect bound at each intermediate time \(s\): \(\|D_v(s)\|\le MCs^2/2\). Its continuous scalar majorant is integrable, so the norm of the position-defect integral is bounded by \(\int_0^tMCs^2/2\,ds=MCt^3/6\). First, the argument verifies that every intermediate time stays in the original horizon, rewrites the scalar coefficient into the majorant's form, applies norm-integral comparison, and uses the elementary integral of a power. The integration variable inside the bound is \(s\), preserving the correct cubic constant. No stationarity or Gaussian moment is assumed implicitly by this deterministic estimate.



For the position defect the first response is bounded by the remaining time: \(\psi_1(h-s)\le h-s\). Thus \(MR\int_0^h(h-s)s\,ds=MRh^3/6\). The extra integration is responsible for the one-higher power of the position error.

\begin{leanbox}{Lean: \texttt{positionDefect\_\allowbreak{}bound}}
\leansource{\texttt{positionDefect\_\allowbreak{}bound} \quad KineticLangevin/\allowbreak{}LocalError.lean, lines 75–95}
\begin{leancode}
\leanline{0}{theorem\ positionDefect\_bound\ (g\ :\ E\ →\ E)\ (M\ :\ ℝ≥0)\ (hg\ :\ LipschitzWith\ M\ g)}
\leanline{4}{(x\ v\ :\ ℝ\ →\ E)\ (γ\ T\ C\ :\ ℝ)\ (hγ\ :\ 0\ ≤\ γ)\ (hC\ :\ 0\ ≤\ C)}
\leanline{4}{(hbound\ :\ ∀\ s\ ∈\ Icc\ 0\ T,\ ‖v\ s‖\ ≤\ C)}
\leanline{4}{(hkin\ :\ ∀\ s\ ∈\ Icc\ 0\ T,\ x\ s\ -\ x\ 0\ =\ ∫\ r\ in\ 0..s,\ v\ r)}
\leanline{4}{\{t\ :\ ℝ\}\ (ht\ :\ t\ ∈\ Icc\ 0\ T)\ :}
\leanline{4}{‖positionDefect\ g\ x\ γ\ t‖\ ≤\ (M\ :\ ℝ)\ *\ C\ *\ t\ \textasciicircum{}\ 3\ /\ 6}
\end{leancode}
\leanlabel{How the Lean proof follows the argument}
\texttt{position\allowbreak{}Defect\_\allowbreak{}bound} gives the formal counterpart of this step. The main intermediate facts are \texttt{hb}, \texttt{hint} in that order. The final step is the algebraic normalization or ordered arithmetic displayed in the proof.
\end{leanbox}

\section{\texorpdfstring{The paper’s three-quarter transformed local constant}{The paper’s three-quarter transformed local constant}}


Put \(B=MCh^2\). The component bounds give \(\|D_v\|\le B/2\) and \(\|D_x\|\le Bh/6\). If \(\gamma h\le1/4\), then \(\|D_v+\gamma D_x\|\le B/2+B/24=13B/24\). The transformed energy at parameter zero is \(\|D_v+\gamma D_x\|^2+\|D_v\|^2\), which is at most \(((13/24)^2+1/4)B^2\le(3B/4)^2\). The proof keeps all quantities nonnegative, squares each norm bound with comparison of nonnegative squares, unfolds the transformed energy and verifies the rational comparison by elementary real inequalities. This derives the constant for the explicit defect integrals; identifying them with an stationary-chain coupling still requires the kinetic process construction and moment estimates.



The transformed norm combines velocity and friction times position. Under the stated step restriction, the cubic position term is bounded by the remaining part of the quadratic velocity-error estimate. The resulting constant must be combined with exact-flow contraction before accumulating errors over many steps.

\begin{leanbox}{Lean: \texttt{localDefect\_\allowbreak{}energy\_\allowbreak{}bound}}
\leansource{\texttt{localDefect\_\allowbreak{}energy\_\allowbreak{}bound} \quad KineticLangevin/\allowbreak{}LocalError.lean, lines 100–133}
\begin{leancode}
\leanline{0}{theorem\ localDefect\_energy\_bound\ \{F\ :\ Type*\}\ [NormedAddCommGroup\ F]\ [InnerProductSpace\ ℝ\ F]}
\leanline{4}{(g\ :\ F\ →\ F)\ (M\ :\ ℝ≥0)\ (hg\ :\ LipschitzWith\ M\ g)}
\leanline{4}{(x\ v\ :\ ℝ\ →\ F)\ (γ\ h\ C\ :\ ℝ)\ (hγ\ :\ 0\ ≤\ γ)\ (hh\ :\ 0\ ≤\ h)\ (hC\ :\ 0\ ≤\ C)}
\leanline{4}{(hsmall\ :\ γ\ *\ h\ ≤\ 1\ /\ 4)}
\leanline{4}{(hbound\ :\ ∀\ s\ ∈\ Icc\ 0\ h,\ ‖v\ s‖\ ≤\ C)}
\leanline{4}{(hkin\ :\ ∀\ s\ ∈\ Icc\ 0\ h,\ x\ s\ -\ x\ 0\ =\ ∫\ r\ in\ 0..s,\ v\ r)\ :}
\leanline{4}{transformEnergy\ γ\ 0\ (velocityDefect\ g\ x\ γ\ h,\ positionDefect\ g\ x\ γ\ h)\ ≤}
\leanline{6}{((3\ /\ 4\ :\ ℝ)\ *\ (M\ :\ ℝ)\ *\ C\ *\ h\ \textasciicircum{}\ 2)\ \textasciicircum{}\ 2}
\end{leancode}
\leanlabel{How the Lean proof follows the argument}
\texttt{local\allowbreak{}Defect\_\allowbreak{}energy\_\allowbreak{}bound} gives the formal counterpart of this step. The main intermediate facts are \texttt{B}, \texttt{h\allowbreak{}B}, \texttt{hv}, \texttt{hp}, \texttt{ht}, \texttt{hsum} in that order. Its remaining equalities are discharged by the simplifications displayed in the proof.
\end{leanbox}

\end{document}
